%% ==============================================================================
%%
%%  Document settings and macro definitions
%%
%% ==============================================================================

\documentclass[a4paper,10pt,bibtotoc]{scrartcl}
\usepackage{a4wide}
\usepackage{usg}
\usepackage{draftcopy}
\usepackage{upgreek}
\usepackage{jsonparse, booktabs}
%\usepackage[landscape]{geometry}
\usepackage{pdflscape}

\newcommand\sumline[2]{%
  \hbox to \linewidth{%
    \strut
    \parbox[t]{.85\linewidth}{\strut #1\strut}\hss
    \parbox[t]{.15\linewidth}{\raggedleft\strut #2\strut}%
  }
}



%% Do NOT remove: required to extract SVN information
\svnInfo $Id$

%% Adjust page footer
\fancyfoot[LE,LO]{LOFAR2-DDD-XXX: Lightning Image Data}
\fancyfoot[RE,RO]{\textsc{lofar} Project}

%% ==============================================================================
%%
%%  Document
%%
%% ==============================================================================

\begin{document}

\title{LOFAR Data Description Document:\\ Lightning Image Data\\
  {\normalsize Document ID: LOFAR2-DDD-010} \\
  {\normalsize Version 0.03.00} \\
   {\normalsize Persistent DOI: XXXXXX} \\
\author{ \parbox{\linewidth}{\centering S.~ter~Veen,  K.~Mulrey, B.H.~Hare H.~Holties} }
\date{\small{Compile Date: \today}} 
}

\maketitle

\section{Table of contents}
\tableofcontents

\listoffigures

\listoftables

\clearpage

%% ______________________________________________________________________________
%%                                                                  Change record




\input common_tables/version_numbering




\clearpage

%%_______________________________________________________________________________
%%                                                                   Introduction

\section{Introduction}
\label{sec:introduction}


\subsection{Purpose and scope}
\label{sec:purpose and scope}

This data product definition document (DDD) describes the internal structure
of and the interface to the LOFAR lightning image data. The data consist of files that contain imaged lighting flashes, produced using TBuf timeseries data and the Lightning Imaging Pipeline. This document does not give details of the pipeline which produces these files.

%R time series data as generated by the Transient Buffers (TBufs). Time series data -- i.e. the digitised voltage output, as received by the individual LOFAR dipoles -- represent the primary input data to the cosmic ray and lightning analysis pipeline(s) and are the most basic form in which the received radio signals are present within the LOFAR
%system. In LOFAR 1.0 these were referred to as the Transient Buffer Boards (TBBs); with the LOFAR 2.0 upgrade, the naming of this functionality changed to Transient Buffer (TBuf) as they do not have their own separate board. This document therefore serves as reference for the implementation of tools that read and write the TBuf data.




\subsection{Applicable Documents}
\input common_tables/applicable_to_all_documents



\begin{center}
  %% Table head
  \tablefirsthead{
    \hline
    \textsc{Reference} & \textsc{Title} & \textsc{Description} \\
    \hline
  }
  \tablehead{
    \multicolumn{3}{r}{\small \textbf{Applicable documents}
      \textsl{continued from previous page}} \\
    \hline
    \textsc{Reference} & \textsc{Title} & \textsc{Description} \\
    \hline
  }
  %% Table tail
  \tabletail{
    \hline
    \multicolumn{3}{r}{\small\sl continued on next page} \\
  }
  \tablelasttail{\hline}
  %% Table contents
  \bottomcaption[List of documents that apply to to this document]{
    List of documents that apply to this Documents.
  \label{tab:applicable list specific}}
  \begin{supertabular}{lp{5.8cm}p{7cm}}
    \shrinkheight{-4.5em}
    \texttt{ICD-001} \cite{lofar.icd.001} & TBB Timeseries-Data & Previous definition of the Transient Buffer Timeseries Dataformat as used by the previous LOFAR system (2009-2022). This document is based on that document. \\
    \hline
  \end{supertabular}
\end{center}


\subsection{Reference applications}
This document serves as reference for the implementation of tools that read and write the lightning data. A non-extensive list is provided here.
\begin{comment}
Add lightning imaging pipeline?
\end{comment}

\begin{center}
  %% Table head
  \tablefirsthead{
    \hline
    \textsc{Name} & \textsc{Title} & \textsc{Purpose} \\
    \hline
  }
  \tablehead{
    \multicolumn{3}{r}{\small \textbf{Applicable documents}
      \textsl{continued from previous page}} \\
    \hline
    \textsc{Reference} & \textsc{Title} & \textsc{Description} \\
    \hline
  }
  %% Table tail
  \tabletail{
    \hline
    \multicolumn{3}{r}{\small\sl continued on next page} \\
  }
  \tablelasttail{\hline}
  %% Table contents
  \bottomcaption[List of tools that read and/or write this data format]{
    List of tools that read and/or write this data format.
  \label{tab:applicable list}}
  \begin{supertabular}{lp{4.8cm}p{6cm}}
    \shrinkheight{-4.5em}
    %\texttt{Cosmic Ray Pipeline} \cite{tools.2} & Cosmic Ray Pipeline & Pipeline to analyse cosmic ray data \\
    \texttt{Lightning mapper} \cite{tools.3} & Lightning imager &  3-D imager for lightning \\
    \hline
  \end{supertabular}
\end{center}
%% ______________________________________________________________________________
%%                                                              Section: Overview

\subsection{Overview}
\label{sec:overview}


This document is structured as follows: %
Section \ref{sec:structure} will describe fundamental overall
structure, including a statement of the primary data product format,
\textcolor{red}{(space separated files with a header)}. These conventions will also include names, meaning, and physical units that may be used to generate and interpret the data files. %
Section \ref{sec:detailed structure} will present a detailed
specification for the data, including a description of the structure
of a LOFAR lightning data naming conventions, units, physical quantities. %

%%_______________________________________________________________________________
%%                                                       Organisation of the data

\clearpage

\section{Organisation of the data}
\label{sec:structure}

\subsection{Conventions}

\input common_tables/types

\input common_tables/notation


\subsection{Data Structure}


Lighting Image Data consists of ``space separated" text files with a header followed by 12 columns of data. 

\subsubsection{Lightning Image Headers}
The header consists of 11 lines, an example of which is shown below.

\begin{verbatim}
v 1
# NOTE: this is data has a very loose cut and the data must be cut further!
# RMS<3,  num_participating_RS>10, and up>0, are reasonable cuts to start with
% Software LOFLI v 2
% antennaset LBA 10-90 
% flashName 19A-1
% timebase_UTC 2019-04-24T19:44:32.504000
! timeID D20190424T194432.504Z
! max_num_data 748792
! data_format i d d d d d d d d d i
unique_ID easting[m] northing[m] up[m] time_from_second[s] latitude longitude
altitude[m] ref_ant_amplitude RMS[ns] location_error_estimate[m] num_participating_RS 

\end{verbatim}

%The lines indicate:
%\begin{itemize}
%    \item version of the data format
%    \item note about data cuts
%    \item note about quality cuts
%    \item software version
%    \item antennaset used in the observation
%    \item lightning flash identifier
%    \item timestamp of the flash
%    \item identifiers and units associated with the data columns
%\end{itemize}



\subsubsection{Lightning Image Source Data}

The following space separated data is organized in rows, each corresponding to a fitted lightning source. Each row consists of:
\begin{itemize}
    \item unique\_ID: identifier for each lightning source
    \item easting [m]
    \item northing[m]
    \item up[m]
    \item time\_from\_second[s]
    \item latitude
    \item longitude
    \item altitude[m]
    \item ref\_ant\_amplitude
    \item RMS[ns] 
    \item location\_error\_estimate[m] 
    \item num\_participating\_RS
\end{itemize}



%A LOFAR TBuf Time-series data file will adhere to the following guidelines:\\

%A LOFAR TBuf Times-series data file will be defined within the context
%f the HDF5 file format.  A LOFAR TBuf Time-series HDF5 %file structure
%will comprise a primary group, a ``root group" in HDF5 %nomenclature,
%which may be considered equivalent to a primary header/data unit (HDU)
%of a standard multi-extension FITS file.  This primary %group will
%consist only of header keywords (“attributes” in HDF5 %nomenclature)
%describing general properties of an observation, along %with pointers
%to contained subgroups.  Those subgroups will comprise %an arbitrary
%number of ``AntennaFieldGroups'' (see %sec~\ref{sec:antenna field group}), where a
%Antenna Field will contain data and meta-data produced %by an
%individual Antenna Field of an indvididual LOFAR %station. Each antenna field group contains one or more %``Dipole Datasets``.

%Figure~\ref{fig:TBB_highlevel} shows the basic organisation of the dataset within
%the HDF5 format. The hierarchical structure essentially follows the
%hierarchical structure of LOFAR itself, i.e. in a top-down approach
%from array through stations through antenna fields down to individual dipoles. The grouping
%of multiple antennas/dipoles within a station is mirrored by the
%collection of dipole datasets into the antenna %field groups.





%%_______________________________________________________________________________
%%                                                    Detailed Data Specification


%%______________________________________________________________________________
%%                                                   Detailed Data Specification

\clearpage

\section{Detailed Data Specification}
\label{sec:detailed structure}

\input common_tables/metadata_intro

\subsection{Naming}
Lightning Impage Data files are named as follows:\\\,\\
%\verb|L<ID>_D<date>T<time>Z_<stationname>_<fieldname>_R<part>_tbb.h5|\\\,\\
%\verb|L<ID>_D<date>T<time>Z_<stationname>_<fieldname>_R<part>_tbb.raw|\\\,\\
%Example:\\
%\verb|L123456789_D20140123T134448.920Z_CS002_LBA_R000_tbb.h5|\\\,%\verb|L123456789_D20140123T134448.920Z_CS002_LBA_R000_tbb.raw|\\\,\\

%In this 
%\begin{itemize}
%    \item L is an identifier for LOFAR. If used for different %instrument, this can change.
%    \item \verb|<ID>| is the subtaskID of the LOFAR TMSS process %managing the data acquisition process. For this version that are the \verb|tbuf_dump| subtasks
%    \item The \verb|<date>| and \verb|<time>| is either the requested start and end time of the data or the actual start and end time of the data. The actual start time of the data may be later then the requested start time and is annotated in the metadata of the file. The time is specified up till millisecond precision. The time is indicated in UTC, as also denoted by the Z.
%    \item The \verb|<stationname>| is the name of the first station in the file.
%    \item \verb|<fieldname>| is the name of the first antenna field in the file. One of (LBA, HBA, HBA0, HBA1, NNF). 
%\end{itemize}
%The lengths of the filename may vary because of the \verb|<ID>| and the \verb|<fieldname>|. The "h5" file contains the header information, as described in this document. The ".raw" file contains the actual data. This split is motivated to easily be able to read the "h5" file, for example to gather feedback. The end user should make sure both the "h5" and the "raw" file are copied over. 


\subsection{The Root Group}
\label{sec:root group}

The LOFAR file hierarchy begins with the top level \verb|`Root`| Group.  This is
the file entry point for the data, and the file node by which navigation of
the data is provided. The \texttt{Root} group will comprise a set of attributes
that describe the underlying file structure, observational metadata, the LOFAR
TBuf data, as well as providing hooks to all groups attached to the \texttt{Root} group.

This section will specify two sets of attributes that will appear in
the \texttt{Root} group: a set of \cla\ (CLA) that will be common to
all LOFAR science data products, and a set of attributes that are
specific to LOFAR TBB Time-series data.  Though these attributes will
all appear together in the \texttt{Root} attribute set, they are
separated in this document in order to demarcate those general LOFAR
attributes that are applicable across all data, and those attributes
that are TBB-specific.\\
In other words,

\begin{verse}
\verb|Root Attributes = Common LOFAR Attributes (CLA) + Supplemental TBB Root Attributes.|
\end{verse}

The \cla\ are the first attributes of any LOFAR file root group.

\subsubsection{\cla}
\label{sec:common lofar attributes}

This section will specify a set of attributes that will be common to
LOFAR science data products.  These ``common LOFAR metadata'' will
appear as attributes at the root level of all LOFAR data files.
\textit{All} LOFAR data products, including TBuf Time-series
\textit{inter alia,} will share a common set of metadata root-level
attributes.  These common LOFAR metadata are to be the first set of attributes of any LOFAR file root group.

Table~\ref{tab:lofar common metadata} lists the {\cla} (CLA) which can
be found in LOFAR observation mode data types. These Attributes are required to be in the Root
Group; if a value is not available for an Attribute, a \verb|`NULL'|
will be used in its place. Special considerations for these keywords: \texttt{OBSERVATION\_ID} contains the subtask ID of the data acquisition specific for this trigger. Other \texttt{OBSERVATION...} keywords consider the data for this trigger/file, for example the \texttt{OBSERVATION\_START\_UTC} is the start of the data dump. The \texttt{OBSERVATION\_FREQUENCY...} keywords refer to the whole band for this filter, or for all filters used. For example the \texttt{110\_190 MHz} filter will provide a min,center,max frequency of 100, 150, 200 Mhz, respectively. \texttt{OBSERVATION\_NOF\_BITS\_PER\_SAMPLE} is 12 for TBuf data. \texttt{FILTER\_SELECTION} and \texttt{ANTENNA\_SET} are usually set to \texttt{NULL}. They are repeated at lower levels. They can be set to the actual value, if and only if a single value is used at lower level within this file, but this is not obligatory as pipelines should better rely on the value per dipole.

\begin{table}[tbp]
  \centering
  \begin{tabular}{|lrp{6cm}|}
    \hline
    \textsc{General LOFAR Group} & \textsc{Value} & \textsc{Description} \\
    \hline
    Root          & \verb|`Root'|    & Top-level LOFAR group type.\\
    \hline
    \hline
    \textsc{TBB Specific Groups} & \textsc{Value} & \textsc{Description} \\
    \hline
    
    Trigger group & \verb|`TriggerGroup'| &  Trigger data group. \\
    Antenna Field group & \verb|`AntennaFieldGroup'|  & Antenna Field data group. \\
    Dipole dataset & \verb|`DipoleDataset'| & Individual Dipole Dataset (Raw Voltage Mode). \\
    % Processing History group &\verb|`ProcessHist'| & This is a Processing History group. \\
    \hline
  \end{tabular}
  \caption{LOFAR TBB Time-series Group Types.}
  \label{tab:group types}
\end{table}


\input common_tables/lofar_common_metadata

%\subsubsection{Additional TBuf Time-series Root Attributes}
%\label{sec:additional root group attributes}

%As explained at the beginning of Sec.~\ref{sec:root group} above, the
%root group of a \textbf{TBuf Time-series data file} will contain a set
%of attributes, which can be broken down into two subsets: 1) a set of
%\cla\ (CLA) that will be common to all LOFAR science data products,
%and 2) a set of attributes that are specific to LOFAR Time-series data
%(see Table~\ref{tab:hdf5 root group}). With the {\cla} already listed
%in section~\ref{sec:common lofar attributes} above, this section will
%focus on the second subset of root group attributes, as they are
%specific to a \textbf{TBuf Time-series data file}.

%\begin{table}[htbp]
%  \centering
%  \begin{tabular}{|llp{9cm}|}
%    \hline
%    \textsc{Field/Keyword} & \textsc{Type} & \textsc{Description} \\
%%%    \hline
%    \hline
%    %\verb|NOF_STATIONS| & \verb|unsigned| &  The number of station groups
%    %available in the file.\\
%    \verb|ANTENNA_FIELD_{AAKKK}_{XBAX}| & Group & Antenna field group %collecting the data
%    from an individual LOFAR station; It consists of the station name with 2
%    characters (upper case) and three digits, followed by the antenna field name (LBA, HBA, HBA0, HBA1.\\
%    \verb|TRIGGER| & Group & Group collecting the
%    parameters associated with the trigger and the output
%    parameters generated by it. \\
%    \hline
%  \end{tabular}
%  \caption{Additional attributes and objects attached to the root
%    group of a TBB time-series data file.}
%  \label{tab:hdf5 root group}
%\end{table}

\clearpage



%%_______________________________________________________________________________
%%                                                    Trigger Group 
\clearpage %% Needed for correct positioning the long trigger group table
           %% w.r.t. other tables.

\paragraph{Trigger Group}

\label{sec:trigger group}

The \verb|TRIGGER| group collects the parameters that caused the TBB to be frozen and the data to be collected. Table~\ref{tab:trigger group} shows the attributes within this group.




\begin{center}
  \tablefirsthead{%
    \hline
    \textsc{Field/Keyword} & \textsc{Type} & \textsc{Software} & \textsc{Description} \\
     && \textsc{Version} &\\
    \hline\hline
  }
  \tablehead{
    \multicolumn{4}{r}{\small\textsl{Table \ref{tab:trigger group}: continued from previous page}}\\
    \hline
    \textsc{Field/Keyword} & \textsc{Type} & \textsc{Software} & \textsc{Description} \\
     && \textsc{Version} &\\
    \hline\hline
  }
  \tabletail{
    \hline
    \multicolumn{4}{r}{\small\textsl{Table \ref{tab:trigger group}: continued on next page}}\\
  }
  \tablelasttail{
    \hline
  }
  \bottomcaption{Attributes attached to a TRIGGER group (TriggerGroup).}
  \begin{supertabular}{|lllp{5.25cm}|}
    \small \texttt{GROUPTYPE} & \verb|string| & DDD v1.0.0 & The value of the grouptype is \verb|`TriggerGroup'|.\\
    \small \texttt{TRIGGER\_TYPE} & \verb|string| & DDD v1.0.0 & Type of trigger. The default value is ``Unknown''.\\
    \small \texttt{TRIGGER\_SUBTYPE} & \verb|string| & DDD v1.0.0 & SubType of trigger. The default value is ``Unknown''.\\
    \small \texttt{TRIGGER\_VERSION} & \verb|string| & DDD v1.0.0 & Version of the trigger algorithm, following Semantic Versioning 2.0.0\footnote{www.semver.org}. The default value is 1.0.0. \\
    \small \texttt{TRIGGER\_ID} & \verb|string| & DDD v1.0.0 & Identifier to match with parameters from the detector/script that sent the trigger. \\
    \small \texttt{TIME} & \verb|unsigned| & DDD v1.0.0 & Unix Timestamp of the trigger in number of seconds since January 1, 1970, 00:00:00 UTC.\\
    \small \texttt{SAMPLE\_NUMBER} & \verb|unsigned| & DDD v1.0.0 & Subsecond timestamp of the trigger expressed as sample number inside the second marked by \verb|TIME|, assuming 5 ns samples. \\
    \small \textit{\texttt{ADDITIONAL\_INFO}} & \verb|string| & DDD v1.0.0 & Free form text field. \\
    \hline
  \end{supertabular}
  \label{tab:trigger group}
\end{center}

The type of trigger is given by the keyword \verb|TRIGGER_TYPE| and \verb|TRIGGER_SUBTYPE|. They contain keywords from Table~\ref{tab:trigger_types}.

\begin{table}[htb]
  \centering
  \begin{tabular}{|lp{9cm}|}
    \hline
    \textsc{Value of} \verb|TRIGGER_TYPE| & \textsc{Description} \\
    \hline
    \hline
    \verb|`Unknown'| & Unknown/unrecognised reason for data return. \\
    \verb|`LORA'| & Trigger to capture a cosmic ray. \verb|TRIGGER_SUBTYPE|  = 'LORA' when triggered from the LORA particle detector.\\
    \verb|`Lightning'| & Triggered in order to capture a lightning strike. \verb|TRIGGER_SUBTYPE| = 'blitzortung' or 'magnetic loop' \\
    \verb|`Manual'| & Manually triggered. \verb|TRIGGER_SUBTYPE| = 'User Portal' or 'test script'\\
   \hline
  \end{tabular}
  \caption{Values of TRIGGER\_TYPE.}
  \label{tab:trigger_types}
\end{table}

Additional keywords for trigger setup or other trigger values can be sent in \verb|ADDITIONAL_INFO|. These are highly detector dependent and will require their own description document. If they are described for a detector, please add the description document in the document list.
\subsection{Header Information}


The first line in the file (and file) is a version line. Which should be a ‘v’ followed by a space, and an integer. e.g.:
v 1,  this document describes version 1. In future, the ‘v’ may be replaced by a ‘b’, which will be a binary version. \textcolor{red}{(not currently defined).}

The header consists of two types of lines, comments and info. 

Comment lines start with a ‘\#’ and the following line is human-readable text, to be ignored by any automated file reading. Comments could include info on who made the file, when, and any special details of the imaging, for example.

Info lines start with a ‘!’ mark, and are followed by a space, then an ID that says what info that line contains. Some types of info are required, some are optional. The order is not guaranteed.

Here are the present types of info lines: 
\begin{enumerate}

\item  (required)

\begin{verbatim}
    ! timeid D20170929T202255.000Z
\end{verbatim}

This gives the time ID of the flash. Starting with D and ending with the milliseconds, including the Z. 
Example:


\item (optional)
\begin{verbatim}
    ID: max\_num\_data
\end{verbatim}

This field fives the maximum number of data rows included. This is for allocating memory efficiently. It should not larger than or equal too, but not too much larger than, the actual number of data points.

\item (required)
\begin{verbatim}
ID: data\_format

\end{verbatim}
This is a space-separated list. Each item is ‘i’ ‘d’ or ‘s’. The number of items is equal to number of columns. Where they indicate the type of that colloid. i.e. ‘i’ means the column is integers, ‘d’ is floating point numbers, and ‘s’ are strings. 
\end{enumerate}

Lines that start with ‘\%’ are machine-readable info lines not specified here. It is not allowed to have a info line (starting with ‘!’) that is not defined below. Lines starting with ‘\%’ allow new types of data-specific info. The first word should be a unique id that indicates the kind of following data. The format of following data is unspecified.

There should always be a space after the \#, \% or !

After the header, the file contains data. Any line that doesn’t start with a !, \#, or \% is data. There should be no lines that start with !, \#, o\%  after the first data line.

In between the header and data there should be a line that describes the columns. The first five columns must be:
\begin{verbatim}
    unique_id distance_east distance_north distance_up time_from_second 
\end{verbatim}

where distance is in m and time is seconds. This descriptive row can contain additional columns, but it must only contain space-separated words. No other punctuation.



%%_______________________________________________________________________________
\subsection{5-d data group with (Time, X, Y, Z, Strength)}
\label{sec:5d}
All following lines must be space-separated data, with same number of columns as the describing row, in the format as specified by the format info line. There should be no quotation marks, nans,  infs, or other punctuation. 
%%                                                    Antenna Field Group 

%\subsection{The Antenna Field Group}
%\label{sec:antenna field group}



%%_______________________________________________________________________________
%%                                                    Dipole Dataset 

%\begin{landscape}
%\subsection{Dipole Dataset}
%\label{sec:dipole dataset}


\section{Data Discovery}
This section describes required/used data for discovery and related information.


Specific parameters for data discovery for the lightning data

%Input brian on Discovery parameters:   number of good sources, duration,  number large flashes,  number small flashes,   list of num good sources per large flash,  list of area per flash,  list of XYZ locations per flash,  list of lat, lons per flash,  list of type per flash (list of strings)

\begin{itemize}
    \item numberOfSources - number of good individual sources of emission - int
    \item duration - duration of the flash, in seconds - float
    \item numberOfLargeFlashes - number of large flashes - int
    \item numberOfSmallFlashes - number of small flashes - int
    \item lightningFlashArea - list of area for lightning flash, maximum radius in meters - (list of float)
    \item lightningLocationXYZ - list of XYZ locations per flash, relative to LOFAR Superterp - alternatively could be split in lightningLocationX , lightningLocationY, lightningLocationZ - (list of float)
    \item lightningLocationLonLat - list of lat,lons per flash in degrees - alternatively could be split in lightningLocationLon , lightningLocationLat - (list of float)
    \item lightningType - list of type per flash (list of strings)
    \item TimeStamp - timestamp of the data. \verb|OBSERVATION_START_UTC| from the 'root' table. - string
    \item triggerType - Type of the trigger, from the trigger table. e.g. Lightning, Cosmic Rays. - string
    \item triggerValue - Value of the trigger, this is the Trigger ID of the Trigger group. - int (?)
\end{itemize}



The following are generic parameters for data discovery
\begin{itemize}
    \item creationDate - datetime - creation date of the object
    \item dataProductIdentifier - identifier of the dataproduct.
    \item dataProductIdentifierName - name of the dataproduct identifier
    \item dataProductIdentifierSource - source of the identifier, typically TMSS. 
    \item dataProductType - Type of the dataproduct. TransientBufferDataProduct.
    \item fileFormat - the format of the file. HDF5
    \item filename - the name of the file
    \item globalname - the full path for storage in the LTA
    \item isDirty - isDirtyFlag. nonZero indicates something is wrong with this data. This is used for indicating system issues. Default false.
    \item isValid - any value greater than 0 means that data is Valid.
    \item processIdentifier - The identifier of the process creating this dataproduct. The subtask ID in TMSS
    \item processIdentifierName - The name of the process creating this dataproduct. The name of the subtask in TMSS.
    \item processIdentifierSource - the source of the process creating this dataproduct. In general, TMSS
    \item releaseDate - The release datetime of the dataproduct. Before the release date, the data is only accesible by the project members, after that, for any LTA user.
    \item releasePolicy - Policy for making the data public.
    \item retentionDate - Datetime until which the data is available in the LTA.
    \item retentionPolicy - Policy for determining the retentionDate
    \item storageWriter - Name of the datawriter that created the file. Default: \verb|tbuf_dumper|.
    \item storageWriterVersion - Version of the storageWriter
\end{itemize}

In addition, the generic observation and projectInformation discovery parameters are also present. 

\begin{comment}
How to deal with LBA + HBA observations in the observation table?
\end{comment}



\subsection{Metadata for Discovery}
The following metadata is used for discovery of data in the archive. %TODO is this not in the table?


%% ------------------------------------------------------------------- References
\section{References}

\subsection{Common LOFAR2 attributes}

\subsection{Naming conventions}

\bibliographystyle{abbrv}
\bibliography{references}


%%_______________________________________________________________________________
%%                                                    Discussion & open questions

\clearpage

\appendix
\section{Appendix}

\subsection{Glossary}
\label{sec:glossary}
\addcontentsline{toc}{section}{\glossaryname}

\input common_tables/lofar_common_glossary


\subsection{Data Types}


\subsection{Notation}


\subsection{FAQ/known issues}


\subsection{Discussion \& open questions}

\subsection{Calibration}
When analyzing the data, a calibration correction is needed to align the different datasets from the separate antennas.

The delay to apply is \verb|AntennaFieldGroup.CLOCK_OFFSET|+\verb|STATION_CALIBRATION.DELAY[DIPOLE_ID]|.

In case only one antenna is enabled for that tile in \verb|TILE ANTENNAS ENABLED| the additional delay for that antenna from the list of \verb|TILE_BEAMFORMER_DELAYS| needs to be added.

\begin{comment}
CHECK SIGNS IN THIS EQUATION
\end{comment}

This delay likely needs to be applied as a phaseshift on the fourier transform of the data.



\subsection{Open questions}

The following table presents an overview of (some of the) known open
questions regarding the format definition:

\begin{center}
  \tablefirsthead{
    \hline
    \textsc{Item} & \textsc{Description} & \textsc{Raised by} \\
    \hline
  }
  \tablehead{
   \multicolumn{3}{r}{\small\textsl{continued from previous page}} \\
    \hline
    \textsc{Item} & \textsc{Description} & \textsc{Raised by} \\
    \hline
  }
  \tabletail{
    \hline
    \multicolumn{3}{r}{\small\textsl{continued on next page}} \\
 }
  \tablelasttail{\hline}
  \begin{supertabular}{lp{12.8cm}l}
    05 & Should PROJECT\_PI name be removed because of GDPR? & S. ter Veen \\
 \end{supertabular}
\end{center}

\subsection{Future enhancements}

Though the file format definition is not intended to undergo
considerable changes once gaining release status, there will be future
enhancements to reflect new insights and address noted short-comings.

\begin{enumerate}
\item Some additional metadata are required when the observation is
  done using the HBA [input thanks to Maaijke Mevius]:
  \begin{enumerate}
  \item At station level the 4x4 numbers which give the relative
    position of the dipoles in a tile, given in HBADeltas.conf. We do
    not really need those for analysis, but if you want to simulate
    the data (or do some sort of selfcal) it might be good to have
    them. \\
    Since these numbers are the same for all tiles within a station, I
    think they can be stored in the station metadata.
  \item The pointing of the tile beam (2 angles) should be added. I
    believe it is possible to have different tiles pointing in
    different directions, thus it should be stored per antenna.
  \item By `tiles', do we mean literal tiles, or tiles/dipoles? [CWJ]
  \end{enumerate}
\end{enumerate}

%%_______________________________________________________________________________
%%                                                                     Appendices

\clearpage





%% --------------------------------------------------------------------- Appendix

\subsection{Requirements}
\label{sec:requirements}


\section{Document history}
\addcontentsline{toc}{section}{Change record}


\subsubsection*{LOFAR-ICD-001 v3.3 $\rightarrow$ LOFAR-DDD-001}

\begin{center}
  %% Table head
  \tablefirsthead{
    \hline
    \textsc{Version} & \textsc{Date} & \textsc{Sections} & \textsc{Description of changes} \\
    \hline
  }
  \tablehead{
    \multicolumn{4}{r}{\small\textsl{continued from previous page}} \\
    \hline
    \textsc{Version} & \textsc{Date} & \textsc{Sections} & \textsc{Description of changes} \\
    \hline
  }
  %% Table tail
  \tabletail{
    \hline
    \multicolumn{4}{r}{\small\textsl{continued on next page}} \\
  }
  \tablelasttail{\hline}
  %% Table contents
  \begin{supertabular}{lllp{9.5cm}}
    4.00.00 & 2025-06-13 & all   & Changes for LOFAR2.0 compared to LOFAR ICD: Removed Subband Groups. Introduced \texttt{AntennaFieldGroup}. Added keywords related to Tile beamformer. Renamed antenna IDs and adjusted appropriate ID keywords. Refined trigger group.   Moved \texttt{DIPOLE\_CALIBRATION} values to a separate group where calibration table can be stored. Adjusted document format to DDD template, including the common metadata group. Moved \texttt{ANTENNA\_SET} and \texttt{FILTER} to new \texttt{AntennaFieldGroup} \\
    4.01.00 & 2018-10-24 & all   & Cleanup of the document. \\
 \end{supertabular}
\end{center}


\subsection{Metadata}
\label{sec:metadata}

Metadata is the auxiliary data stored along with the time series data need to
provide all the necessary information for automated processing of the data.
This data can either be stored directly in the data set or it can be stored
in an external database. In the latter case the data set must contain a
pointer to the correct entry in that database (e.g. the antenna-id is needed
to get the antenna position).

\begin{itemize} \parskip 0pt
\item DAQ mode, including Samplerate, Filters, etc.
\item Timing information:
  \begin{itemize}
  \item Trigger time relative to recorded data segement
  \item Timing of the data streams relative to the trigger
  \item Timing of the data streams relative to each other with sub-sample
    accuracy; This can be implicit, e.g. all data streams of a station start
    at the same time. Fields that can be in stored in an external database.
  \end{itemize}
\item List of RFI sources identified by the station calibration, including
  \textbf{direction}, \textbf{center frequency} and \textbf{peak strength}.
  \begin{list}{$\rightarrow$}{}
  \item What does this actually mean: the properties of the single channel
    containing the highest signal level or the parameters obtained from fitting
    e.g. a Gaussian to a segment of the spectrum?
  \end{list}
\item Dispersion measure of the ionosphere (at this point in time and space)
\item Health information about the antennas
\end{itemize}

\subsubsection{System monitoring and system health.}

Information on the status of the various (hardware) components at a LOFAR station
will be stored inside the PVSS database; a description of the datapoint-types and
datapoints can be found in the \verb|MAC/Deployment/data/PVSS| branch of the
LOFAR code repository (see Table~\ref{tab:PVSS} for an excerpt).

\begin{table}[htb]
  \centering
  \begin{tabular}{|lll|}
    \hline
    \textsc{Database entry} & \textsc{Field} & \textsc{Format} \\
    \hline
    \hline
    CalCtlr & connected    & unsigned int \\
            & obsname      & string \\
            & antennaArray & string \\
            & filter       & string \\
            & nyquistzone  & int \\
            & rcus         & string \\
    \hline
    ObservationControl & claimPeriod & int \\
    & preparePeriod   & int \\
    & startTime       & string \\
    & stopTime        & string \\
    & subbandList     & string \\
    & beamletList     & string \\
    & bandFilter      & string \\
    & nyquistzone     & int \\
    & antenneArray    & string \\
    & receiverList    & string \\
    & sampleClock     & int \\
    & measurementSet  & string \\
    & stationList     & string \\
    & inputNodeList   & string \\
    & BGLNodeList     & string \\
    & storageNodeList & string \\
    \hline
  \end{tabular}
  \caption[Excerpt from the list of entries into the PVSS database.]{
    Excerpt from the list of entries into the PVSS database. The
    definitions of the datapointtypes and datapoints can be found in the
    \texttt{MAC/Development/data/PVSS} branch of the LOFAR code repository.}
  \label{tab:PVSS}
\end{table}

In order to later store certain system health information along with the other
data, parameters need to be subscribed to at the definition of the observation.

\subsection{Visualization}

Past experience with the software for the LOPES experiment has shown, that is
is crucial to provide the user with a variety of way to graphically inspect
the data. This not only includes visualization of the time-series/FFT/etc.
data themselves, but also displaying the various data with the wider context
of the experimental setup (e.g. geographical distribution of the antennas
w.r.t. to the particle detector setup)

\begin{itemize}
\item Display of the standard data products (also see documentation of the
  \lopestools\ \verb|DataReader|): ADC, Voltage, FFT, Calibrated FFT,
  RFI-filtered FFT, Cross-Corr. Spectra, Visibilities
\item Display of the (intermediate) data products generated from the input
  data, e.g. dynamic spectra, multi-dimensional skymaps
\item Flags and weights associated with the data, e.g. filter curves, antenna
  gain curves, etc.
\item Station layout, i.e. positions of the (selected/excluded) antennas
\item antenna power level distribution over the area of the station/array
  (this is very similar to the type of event display as known from particle
  physics experiments)
\item mapping of the (local) RFI via an (Azimuth,Frequency) plot centered on
  the position of a certain station
\item geographical locations of identified RFI sources; such a plot should
  also indicate the  frequency band of the RFI (via label or bar etc.)
\item geographical distribution/location of antennas which generated a trigger
  signal, failed, etc.
  \begin{itemize}
  \item[$\rightarrow$] VR setting, combining geographical information (e.g.
    map of the Netherlands) with the location of localized CR-/TS-events
  \item[$\Rightarrow$] in a cave-like setup we actually can perform a
  fly-through, which would be ideal for outreach purposes!
  \end{itemize}
\item cummulative geographical distribution of detected CR events
\item total power per LOFAR station (geographically distributed)
\item overlay of CR data with information from other sensors (e.g. weather
  radar images, temperatures, etc.)
\end{itemize}
A number of the before-mentioned displayes should be interactive, in the sense
that the user should be able to perform data selection from the graphical
display (e.g. by drawing a circle around the core of the CR air shower,
thereby selecting the antennas included in the data analysis step).





\end{document}
